% Part:sets-functions-relations
% Chapter: size-of-sets
% Section: non-enumerability

\documentclass[../../../include/open-logic-section]{subfiles}

\begin{document}

\olfileid[ar]{sfr}{siz}{nen}

\olsection{مجموعات \printtoken{S}{nonenumerable}}

\begin{editorial}
المقصود هنا إثبات أن $\Bin^\omega$ و$\Pow{\PosInt}$ غير قابلتين للتعداد، على التعريف المتقدم في \olref[enm]{sec}. وروعي في هذا العرض أن يكون أيسر مدخلًا وأبسط بيانًا، مع زيادة شيء من التفصيل، من العرض الذي في \olref[enm-alt]{sec}.
\end{editorial}

من المجموعات ما لا نهاية له، كمجموعة الأعداد الصحيحة الموجبة~$\PosInt$. وقد كانت المجموعات اللانهائية التي تقدمت أمثلتها كلها !!{enumerable}؛ وليس ذلك شأن كل مجموعة لا نهائية. فالمجموعة التي تنتفي عنها هذه الخاصية هي \emph{!!{nonenumerable}}.

وقد يستغرب المرء أول الأمر وجود مجموعات !!{nonenumerable}. ووجه ذلك أن المجموعة غير الخالية !!{enumerable}~$A$ تقبل دالة~!!a{surjective} $f \colon \PosInt \to A$، وأما المجموعة !!{nonenumerable} فلا تقبلها. فمهما أرسلنا !!{element}s~$\PosInt$، على ما فيها من اللانهاية، إلى~$A$، لم نستوف بذلك~$A$. فكأن !!{element}s~$A$ «أكثر عددًا» حتى من الأعداد الصحيحة الموجبة التي لا نهاية لها.

\textbf{تنبيه تصحيحي على الأصل.} قُيّدت دعوى وجود الدالة الشاملة بكون
المجموعة غير خالية؛ إذ لا توجد دالة كلية من الصحيحات الموجبة إلى المجموعة الخالية،
مع أن الخالية داخلة في المجموعات المعدودة بحسب التعريف المتقدم.

وطريق إثبات أن المجموعة غير الخالية !!{nonenumerable} أن ننفي إمكان تلك الدالة الشاملة؛ ومعناه أن عناصر~$A$ لا تستوعبها قائمة لا نهائية تمتد في جهة واحدة. ويكفينا لذلك أن نبين أن كل قائمة من !!{element}s~$A$ يفوتها عنصر على الأقل، أي إن الدالة $f\colon \PosInt \to A$ لا تكون !!{surjective} أيًّا كانت. ولهذا نستعمل \emph{طريقة كانتور القطرية}: نأخذ قائمة من !!{element}s~$A$، ولتكن $x_1$، $x_2$، \dots، ثم نبني عنصرًا من~$A$ يقتضي نفس بنائه خروجه عن القائمة.

ولنبدأ بالمجموعة~$\Bin^\omega$: عناصرها جميع المتتاليات اللانهائية التي لا فجوات فيها، ولا تشتمل إلا على~$0$ و~$1$.

\begin{thm}
\ollabel{thm:nonenum-bin-omega}
$\Bin^\omega$~!!{nonenumerable}.
\end{thm}

\begin{proof}
لنفرض خلاف المطلوب، فتكون $\Bin^\omega$ !!{enumerable}، وتوجد قائمة $s_{1}$، $s_{2}$، $s_{3}$، $s_{4}$، \dots{} تستوفي !!{element}s~$\Bin^\omega$. وكل~$s_i$ في هذه القائمة متتالية لا نهائية من~$0$ و~$1$. فإذا كان~$j$ رتبة العنصر، و~$i$ رتبة المتتالية في القائمة، سمينا ذلك العنصر $s_i(j)$. فالمتتالية ذات الرتبة~$i$، وهي $s_i$، صورتها
\[
s_i(1), s_i(2), s_i(3), \dots
\]

ونعرض القائمة في مصفوفة، نجعل لكل متتالية~$s_i$ صفًّا ترتب فيه عناصرها:
\[
\begin{array}{c|c|c|c|c|c}
& 1 & 2 & 3 & 4 & \dots \\\hline
1 & \mathbf{s_{1}(1)} & s_{1}(2) & s_{1}(3) & s_1(4) & \dots \\\hline
2 & s_{2}(1)& \mathbf{s_{2}(2)} & s_2(3) & s_2(4) & \dots \\\hline
3 & s_{3}(1)& s_{3}(2) & \mathbf{s_3(3)} & s_3(4) & \dots \\\hline
4 & s_{4}(1)& s_{4}(2) & s_4(3) & \mathbf{s_4(4)} & \dots \\\hline
\vdots & \vdots & \vdots & \vdots & \vdots & \mathbf{\ddots}
\end{array}
\]
فالأعداد المحاذية للصفوف تبين رتب المتتاليات في القائمة $s_1$، $s_2$، \dots؛ والأعداد فوق الأعمدة تبين رتب !!{element}s كل متتالية على حدتها. فالرمز $s_{1}(1)$، مثلًا، اسم للعدد الواقع أولًا، سواء كان $0$ أم~$1$؛ فهو أول !!{element} في المتتالية $s_{1}$، وعلى هذا القياس سائر الخانات.

ونستخرج من هذه القائمة متتالية لا نهائية $\overline{s}$، مؤلفة من~$0$ و~$1$، لا تكون واحدًا من أعضائها. وإنما يتبع تعريف $\overline{s}$ القائمة المعطاة $s_1$، $s_2$، \dots. فأي قائمة لا نهائية من متتاليات لا نهائية من~$0$ و~$1$ تمدنا بمتتالية~$\overline{s}$ لا تظهر فيها.

وتعيين $\overline{s}$ إنما يكون بتعيين !!{element}sها كلها؛ أي بتعيين $\overline{s}(n)$ لكل $n \in \PosInt$. فنقرأ قطر المصفوفة المتقدمة---ومن هنا جاءت تسمية الطريقة---ونبدل كل $1$ بـ~$0$، وكل $0$ بـ~$1$. وحاصل هذا أن قيمة $\overline{s}(n)$ تكون $0$ أو $1$ تبعًا لكون !!{element} ذي الرتبة~$n$ من القطر، أعني $s_n(n)$، هو $1$ أو $0$:
\[
\overline{s}(n) =
\begin{cases}
1 & \text{إذا كان $s_{n}(n) = 0$}\\
0 & \text{إذا كان $s_{n}(n) = 1$}.
\end{cases}
\]
ويقوم مقام هذا التعريف بالحالات الحد الموجز
$\overline{s}(n) = 1 - s_n(n)$.

وقد تحقق أن $\overline{s}$ متتالية لا نهائية من~$0$ و~$1$؛ إذ هي مقابلة المتتالية التي نقرؤها من قيم~$0$ و~$1$ على القطر، بعد قلب كل قيمة. فهي إذن، أي $\overline{s}$، !!a{element} من~$\Bin^\omega$. وأما امتناع ورودها في القائمة $s_1$، $s_2$، \dots{} فبيانه الآتي.

أما أول متتالية، $s_1$، فلا تساويها؛ لأنها تخالف $s_1$ في أول !!{element}. ومهما تكن $s_1(1)$ فقد جعلنا $\overline{s}(1)$ مقابلها. وأما الثانية فلا تساويها أيضًا، إذ تخالف $\overline{s}$ المتتالية $s_2$ في موضعها الثاني: إن كانت $s_2(2)$ هي $0$، كانت $\overline{s}(2)$ هي $1$، وبالعكس. وهكذا يجري الأمر في سائر المواضع.

ولتقرير ذلك على العموم، لو وردت $\overline{s}$ في القائمة لكان لها رتبة~$k$، ولكان $\overline{s} = s_{k}$. وتساوي المتتاليتين هو اتفاقهما في كل موضع؛ فيلزم لكل~$n$ أن $\overline{s}(n) =
s_{k}(n)$. فإذا خصصنا هذا الحكم بالحالة $n = k$ لزم $\overline{s}(k) = s_{k}(k)$. ولكن $s_k(k)$ لا تخرج عن $0$ و~$1$. فإن كانت $0$، كانت $\overline{s}(k)$ هي~$1$ بمقتضى تعريف $\overline{s}$. وإن كانت $s_k(k) = 1$ اقتضى تعريف $\overline{s}$ أن $\overline{s}(k) = 0$. فالحالتان تفضيان إلى $\overline{s}(k) \neq s_{k}(k)$، وهو التناقض.

فقد أخذنا قائمة من !!{element}s~$\Bin^\omega$ هي $s_1$، $s_2$، \dots{}، ومنها بنينا المتتالية~$\overline{s}$ وبرهنا أنها لا ترد فيها. وهي مع ذلك متتالية من~$0$ و~$1$، ما دامت المتتاليات~$s_i$ كلها من~$0$ و~$1$؛ أي إن $\overline{s} \in
\Bin^\omega$. فلا تستوفي قائمة \emph{جميع} !!{element}s~$\Bin^\omega$، لأن أي قائمة مفروضة تتيح بناء متتالية~$\overline{s}$ خارجة عنها. وهكذا يبطل فرض القائمة المستوفية لجميع متتاليات~$\Bin^\omega$.
\end{proof}

\begin{explain}
إنما سمي هذا الاستدلال قطريًّا لأخذ تعريف~$\overline{s}$ من قطر المصفوفة. وليس وجود المصفوفة شرطًا في الاستدلال القطري؛ فقد نستعمل معناه لإثبات أن مجموعات ليست !!{enumerable}، من غير أن نرسم مصفوفة أو نعين قطرًا ظاهرًا.
\end{explain}

\begin{thm}
\ollabel{thm:nonenum-pownat}
$\Pow{\PosInt}$ ليست !!{enumerable}.
\end{thm}

\begin{proof}
نستعمل المعنى نفسه: فكل قائمة من مجموعات جزئية من~$\PosInt$ تفوتها مجموعة جزئية من $\PosInt$. ولتكن قائمة المجموعات الجزئية من~$\PosInt$ المعطاة
\[
Z_{1}, Z_{2}, Z_{3}, \dots
\]
ونضع مجموعة~$\overline{Z}$ بهذه الصفة: لكل $n \in \PosInt$، يكون $n \in \overline{Z}$ إذا وفقط إذا كان $n \notin Z_{n}$، أي
\[
\overline{Z} = \Setabs{n \in \PosInt}{n \notin Z_n}
\]
فالمجموعة $\overline{Z}$ لا تضم إلا أعدادًا صحيحة موجبة، شأن كل~$Z_n$ بحسب الفرض؛ ولذلك $\overline{Z} \in
\Pow{\PosInt}$. وأما $\overline{Z}$ نفسها فخارجة عن القائمة؛ ويكفي لإثبات ذلك أن نبرهن، لكل $k \in \PosInt$، على $\overline{Z} \neq
Z_k$.

خذ $k \in \PosInt$ كيف اتفق. إن تعريف $\overline{Z}$ يقضي، لكل $n \in \PosInt$، بأن $n \in \overline{Z}$ يكافئ $n \notin Z_n$. فبوضع $n=k$ نجد أن $k \in \overline{Z}$ يكافئ $k \notin Z_k$. وهذا وحده يثبت $\overline{Z} \neq Z_k$؛ إذ $k$ !!a{element} لإحداهما دون الأخرى، فليست $\overline{Z}$ و$Z_k$ متفقتين في !!{element}sهما. وحيث لم نخص $k$ بشيء، لم تكن $\overline{Z}$ واحدة من $Z_1$، $Z_2$، \dots.
\end{proof}

\begin{explain}
وإن لم نذكر القطر في البرهان، أمكن إظهاره بتصوير المجموعات $Z_1$، $Z_2$، \dots{} صفوفًا في مصفوفة، يوضع فيها كل !!{element}~$j \in Z_i$ في العمود ذي الرتبة~$j$. فلو كانت المجموعات الأربع الأولى $\{1,2,3,\dots\}$، $\{2, 4, 6, \dots\}$، $\{1,2,5\}$، و$\{3,4,5,\dots\}$، لكان أول المصفوفة
\[
\begin{array}{r@{}rrrrrrr}
  Z_1 = \{ & \mathbf{1}, & 2, & 3, & 4, & 5, & 6, & \dots\}\\
  Z_2 = \{ &  & \mathbf{2}, &  & 4, &  & 6, & \dots\}\\
  Z_3 = \{ & 1, & 2, &  &  & 5\phantom{,} &  & \}\\
  Z_4 = \{ &  &  & 3, & \mathbf{4}, & 5, & 6, & \dots\}\\
  \vdots & & & & & \ddots
\end{array}
\]
فلاستخراج $\overline{Z}$ ننزل على القطر: كل عدد وجدناه عليه حذفناه، وكل قيمة~$j$ كانت خانتها القطرية في الصف والعمود ذوي الرتبة~$j$ خالية أدرجناها. ففي هذا المثال نحذف $1$ و$2$، ونثبت~$3$، ونحذف~$4$، ونمضي على هذا الوجه.
\end{explain}

\begin{prob}
أقم حجة قطرية تبين أن $\Pow{\Nat}$ !!{nonenumerable}.
\end{prob}

\begin{prob}\label{sfr:siz:nen:prob:f-posint}
بين باستدلال قطري صريح أن مجموعة الدوال $f \colon \PosInt \to \PosInt$ !!{nonenumerable}. والمطلوب أنه متى أعطيت قائمة $f_1$، $f_2$، \dots{} من الدوال، لكل واحدة منها النوع $f_i\colon
\PosInt \to \PosInt$، أمكن إيجاد دالة $\overline{f}\colon \PosInt \to
\PosInt$ لا ترد في القائمة.
\end{prob}

\end{document}
